\newcommand{\auditoriaid}{A02 -- Antecedentes y literatura TOE}
% Preambulo comun de las auditorias del informe E1 (revision_e1/)
\documentclass[11pt,a4paper]{article}
\usepackage[utf8]{inputenc}
\usepackage[T1]{fontenc}
\usepackage[spanish,es-noshorthands]{babel}
\usepackage{lmodern}
\usepackage{microtype}
% Sin patrones de guionado en espanol en esta instalacion (ver A10): se
% desactiva el guionado automatico para evitar cortes con reglas del ingles.
\hyphenpenalty=10000
\exhyphenpenalty=10000
\usepackage[a4paper,margin=2.2cm]{geometry}
\usepackage{booktabs,array,longtable}
\usepackage{graphicx,caption,subcaption}
\usepackage{xcolor}
\usepackage{enumitem}
\usepackage{fancyhdr}
\usepackage{amsmath}
\usepackage{tikz}
\usetikzlibrary{arrows.meta,positioning}
\usepackage{natbib}
\usepackage{xurl}
\usepackage[hidelinks]{hyperref}
\urlstyle{tt}
\newcommand{\archivo}[1]{\url{#1}}
\newcolumntype{L}[1]{>{\raggedright\arraybackslash}p{#1}}
\definecolor{alta}{RGB}{180,30,30}
\definecolor{media}{RGB}{200,120,0}
\definecolor{baja}{RGB}{60,110,160}
\definecolor{cajafondo}{RGB}{245,247,250}
\newcommand{\sevA}{\textcolor{alta}{\textbf{Alta}}}
\newcommand{\sevM}{\textcolor{media}{\textbf{Media}}}
\newcommand{\sevB}{\textcolor{baja}{\textbf{Baja}}}
\setlength{\parskip}{0.35em}
\setlength{\emergencystretch}{3em}
\setlength{\parindent}{0pt}
\pagestyle{fancy}
\fancyhf{}
\fancyhead[L]{\small Auditoría del Informe del Entregable~1 -- \auditoriaid}
\fancyhead[R]{\small \thepage}
\renewcommand{\headrulewidth}{0.4pt}
% Entorno de hallazgos: ID | Severidad | Ubicacion | Hallazgo y evidencia | Correccion
\newenvironment{hallazgos}{%
\small
\begin{longtable}{@{}L{1.15cm}L{1.25cm}L{1.9cm}L{6.0cm}L{\dimexpr\textwidth-10.3cm-10.6\tabcolsep\relax}@{}}
\toprule
\textbf{ID} & \textbf{Sev.} & \textbf{Ubicación} & \textbf{Hallazgo y evidencia} & \textbf{Corrección aplicada} \\
\midrule\endfirsthead
\toprule
\textbf{ID} & \textbf{Sev.} & \textbf{Ubicación} & \textbf{Hallazgo y evidencia} & \textbf{Corrección aplicada} \\
\midrule\endhead
\bottomrule\endfoot
}{\end{longtable}}
% Macros compartidas con el informe revisado (las secciones propuestas las usan)
% Macros compartidas entre el informe revisado y las auditorias.
% Cifras verificadas contra los archivos del repositorio (ver A01/A04/A07):
\newcommand{\nUniverso}{102}   % source_id CMIP6 con tos/Omon en ESGF (00b, model_availability_priority.csv)
\newcommand{\nSeed}{41}        % publican los cuatro experimentos (config/models_seed_cmip6.csv)
\newcommand{\nSel}{40}         % completos + QC PASS (models_inventory_final.csv)
\newcommand{\nComb}{160}       % combinaciones modelo-experimento PASS (qc_report.csv)
\newcommand{\nNoEnc}{44}       % config/models_no_encontrado_44.csv
\newcommand{\nParcial}{18}     % 102 - 40 - 44
% Fecha de la portada: fija (no \today), editable en un solo lugar.
\newcommand{\fechaentrega}{28 de septiembre de 2026}
% Estado de codigo que documenta el E1 (antes de los cambios del E2)
\newcommand{\commitEone}{\texttt{8ee6184}}

% En las auditorias el texto propuesto se muestra fuera del informe: las
% referencias cruzadas se imprimen con el nombre de la seccion destino.
\makeatletter
\def\etq#1#2{\expandafter\def\csname etq@#1\endcsname{#2}}
\etq{sec:resumen}{Resumen}
\etq{sec:alcance}{Alcance}
\etq{sec:marco}{Revisión científica}
\etq{sec:antecedentes}{Antecedentes}
\etq{sec:referencias}{Literatura sobre TOE}
\etq{tab:referencias}{bibliografía TOE}
\etq{sec:similitudes}{Niño~3.4 frente a Niño~1+2}
\etq{sec:toe_metodos}{Elección de los métodos de TOE}
\etq{tab:metodos_toe}{métodos TOE}
\etq{eq:szabo_signal}{ecuación señal M1}
\etq{eq:szabo_V}{ecuación V(t)}
\etq{eq:szabo_T}{ecuación T(t)}
\etq{eq:szabo_toe}{ecuación TOE M1}
\etq{eq:tod}{ecuación ToD}
\etq{sec:datos}{Datos}
\etq{fig:cmip6_contribution}{contribución institucional}
\etq{sec:tabla_modelos}{Modelos seleccionados}
\etq{tab:modelos}{modelos}
\etq{sec:desviacion}{Selección de modelos}
\etq{sec:dominios_periodos}{Dominios y periodos}
\etq{fig:region_nino}{regiones Niño}
\etq{sec:arquitectura}{Arquitectura del pipeline}
\etq{fig:flujograma}{Flujograma}
\etq{sec:scripts}{Descripción de los pasos}
\etq{sec:paso00b}{Paso 00b}
\etq{sec:paso01}{Paso 01}
\etq{sec:paso02}{Paso 02}
\etq{sec:paso04}{Paso 04}
\etq{sec:paso05}{Paso 05}
\etq{sec:paso06}{Paso 06}
\etq{sec:paso07}{Paso 07}
\etq{sec:reporte}{Reporte de estado}
\etq{sec:resultados}{Resultados}
\etq{fig:qc}{matriz de estado}
\etq{tab:cuentas}{conteo de modelos}
\etq{fig:mapas2d}{mapas 2D}
\etq{fig:series}{series}
\etq{fig:boxplot}{boxplots}
\etq{sec:roadmap}{Próximos pasos}
\etq{tab:roadmap}{próximos pasos}
\etq{sec:conclusiones}{Conclusiones}
\etq{sec:recomendaciones}{Recomendaciones}
\makeatother

\makeatletter
\AtBeginDocument{\renewcommand{\ref}[1]{\@ifundefined{etq@#1}{\textit{[#1]}}{\textit{\csname etq@#1\endcsname}}}}
\makeatother
% Texto propuesto: secciones degradadas un nivel y en caja
\newenvironment{propuesto}{%
  \par\medskip\noindent\textcolor{baja}{\rule{\textwidth}{0.8pt}}\par
  \noindent{\small\textbf{Inicio del texto propuesto para el informe revisado}}\par\medskip
  \begingroup
  \let\section\subsection\let\subsection\subsubsection\let\subsubsection\paragraph
}{%
  \endgroup
  \par\noindent{\small\textbf{Fin del texto propuesto}}\par
  \noindent\textcolor{baja}{\rule{\textwidth}{0.8pt}}\par\medskip}
% Recuadro de actualizacion posterior (decisiones del autor)
\newcommand{\actualizacion}[1]{\par\medskip\noindent\fcolorbox{media}{media!8}{\parbox{\dimexpr\linewidth-2\fboxsep-2\fboxrule\relax}{\small\textbf{Actualización del 29-09-2026 (decisiones del autor).} #1}}\par\medskip}

\begin{document}

\begin{center}
{\Large\bfseries Auditoría A02: Antecedentes y literatura sobre TOE}\\[0.4em]
{\normalsize Informe del Entregable~1, Secciones 3.1 y 3.2 (\texttt{informe\_e1\_smnv3.tex}, commit \texttt{330fd2e})}\\[0.2em]
{\small Revisión: \fechaentrega}
\end{center}

\section{Objeto}

Se auditan los Antecedentes (Sección~3.1) y la Literatura sobre TOE (Sección~3.2, con su cuadro de bibliografía). Cada afirmación atribuida a un autor se contrastó contra el texto completo de la fuente: la colección de PDF del proyecto (\archivo{~/personal/loca2026/smn/files_pdf/}), convertida a texto para buscar los pasajes, y la API de Crossref para los metadatos. No se aceptó ninguna afirmación verificada solo contra un resumen cuando el texto completo estaba disponible.

\section{Fuentes contrastadas}

\begin{center}\small
\begin{tabular}{@{}L{2.6cm}L{4.2cm}L{\dimexpr\textwidth-6.8cm-4\tabcolsep\relax}@{}}
\toprule
Cita & Archivo consultado & Pasajes clave verificados \\
\midrule
hawkins2009 & \archivo{hawkins2009narrow.pdf} & Apéndice: polinomio de 4.º orden 1950--2099; Ec.~(7) $T(t)=V+S(t)+M(t)$, fuentes independientes; medias decadales; ponderación de modelos. \\
hawkins2012 & \archivo{hawkins2012toe.pdf} & Señal: regresión de la temperatura local sobre la global suavizada con polinomio de 4.º orden; ruido: desviación estándar interanual; umbrales S/N $>1$ y $>2$; incertidumbre del TOE de 30 a 60 años. \\
szabo2016 & \archivo{zsoba2014uncert.pdf} (capítulo 12) & Ecs.~12.1--12.9; $V(t)$ recalculada cada año sobre los 30 años previos y suavizada con media móvil de 10 años; medias anuales y estacionales; sin ponderación; sin test de significancia; TOE con $\sqrt{V}$ y umbral ``usualmente $+1$ o $-1$''. \\
bruno2023 & \archivo{ramirez2023uncer.pdf} & Portada: \emph{Trabajo de Suficiencia Profesional}; ``cincuenta y tres'' modelos; 1981--2010; SSP2-4.5/5-8.5; Pearson, BIAS, RMSE frente a ERA5; diagramas de Taylor; Anexo~6 con V, S/N y TOE. \\
alvarez2024 & \archivo{alvarez2024tos.pdf} & Portada: \emph{Trabajo de Suficiencia Profesional}; TSM y nivel del mar; costas de Sudamérica y Pacífico tropical; CMIP5. \\
takahashi2019 & \archivo{takahashi2017enso.pdf} & Resumen: FEP cálido y Pacífico central frío; ondas Kelvin con ``little role'' en el inicio; vientos del norte, ZCIT. \\
Resto del cuadro & PDF homónimos y Crossref & Ver hallazgos A02-10 a A02-14. \\
\bottomrule
\end{tabular}
\end{center}

\section{Hallazgos}

\begin{hallazgos}
A02-1 & \sevA & 3.1, Hawkins y Sutton & ``\citet{hawkins2012} formalizó el TOE \ldots\ entendiendo la señal como el ajuste polinomial de la variable de interés''. En HS2012 la señal local es la regresión de la temperatura local sobre la temperatura \emph{global} suavizada (polinomio de 4.º orden). El ajuste polinomial directo de la variable es de HS2009 y de Szabó. & Se describen por separado HS2009/2011 (partición, polinomio, $T=V+M+S$) y HS2012 (regresión sobre la global, umbrales 1 y 2, incertidumbre del TOE). \\
A02-2 & \sevA & 3.1, Szabó & Atribuye a Szabó ``tres modificaciones'', entre ellas que ``las tres fuentes se combinan de forma aditiva''. La suma de varianzas ya está en HS2009, Ec.~(7). Las modificaciones reales, según el propio capítulo, son: $V(t)$ variable (30 años previos), medias anuales/estacionales en vez de decadales, sin ponderación de modelos, y TOE con $V$ como único ruido. & Se listan las cuatro modificaciones documentadas. \\
A02-3 & \sevA & 3.2, párrafo final & ``la elección metodológica desplaza la fecha resultante en décadas (comparar \ldots\ la partición aditiva de Szabó frente a la raíz de suma de cuadrados usada en otros trabajos)''. Sumar varianzas y combinar desviaciones estándar en raíz de la suma de cuadrados dan la misma magnitud, $\sqrt{V+M+S}$: no es una fuente de discrepancia. & Se reemplaza por ejemplos verificados: umbrales distintos (1; 1 y 2; 1{,}6; 2), incertidumbre del TOE de 30 a 60 años (HS2012) y diferencias de más de 40 años entre percentiles (Chadwick). \\
A02-4 & \sevM & 3.1, Bruno & ``54 modelos evaluados''. La tesis dice ``cincuenta y tres''; su Tabla~1 numera hasta 54 porque MPI-ESM-1-2-HAM ocupa una celda combinada (filas 43--44). Se verificaron 53 nombres distintos. & 53 modelos. \\
A02-5 & \sevM & 3.1 y bibliografía & Bruno y Álvarez son trabajos de suficiencia profesional, pero el encabezado dice ``Tesis de Bruno'' y la bibliografía (\texttt{@mastersthesis}) imprime ``Master's thesis''. & Encabezado ``Trabajo de Bruno Ramírez''; se agrega \texttt{type = \{Trabajo de suficiencia profesional\}} en la bibliografía revisada. \\
A02-6 & \sevM & 3.2, introducción & ``incluyendo los antecedentes señalados por el TdR''. El TdR no nombra antecedentes ni autores (ver A01-4). & ``los antecedentes de la Sección 3.1''. \\
A02-7 & \sevM & 3.1, Szabó & El texto incluye el DOI (``adopta \ldots\ (DOI 10.1007/\ldots)'') como si el DOI fuera un documento del TdR. Ya figura en la bibliografía. & Se quita del cuerpo. \\
A02-8 & \sevM & 3.1, Takahashi & ``no necesariamente por ondas Kelvin'' atenúa lo que el artículo afirma (``had little role in its initiation'') y omite el mecanismo principal (vientos del norte a través del ecuador, ZCIT al sur del ecuador). No menciona la conclusión central: dos tipos de eventos de fuerte impacto en el Pacífico oriental lejano. & Se precisa el mecanismo y se agrega la conclusión, que justifica tratar Niño~1+2 por separado. \\
A02-9 & \sevB & 3.1 & Comentario \texttt{\% REVISAR} sobre la fórmula repetida ``es el antecedente más directo'' en Bruno y Álvarez. & Se quita la autocalificación en Álvarez; se conserva una sola en Bruno, justificada. \\
A02-10 & \sevM & Cuadro, Giorgi & ``Origen de la formalización S/N; precedente de no emergencia amazónica''. Giorgi y Bi (2009) es contemporáneo de HS2009, no su origen, y para la cuenca amazónica reportan emergencia ``in the late decades or beyond'', no ausencia de emergencia. & ``TOE de \emph{hot-spots}; emergencia tardía en la cuenca amazónica''. \\
A02-11 & \sevM & Cuadro, Chadwick & ``Chile central'', ``TOE hidroclimático probabilístico''. Las cuencas son Limarí (semiárida), Maipo y Maule; el título dice ``Semiarid and Mediterranean Chile''; el método se basa en potencia estadística. & ``Chile semiárido y mediterráneo''; ``TOE local por potencia estadística''. \\
A02-12 & \sevM & Cuadro, Larson & ``Océano global, fronteras orientales''. Omite que el análisis es sobre el periodo histórico (CESM2, 1900--2014), un dato que cambia la lectura de la Sección~3.3 (ver A03). & ``Océano global (CESM2, 1900--2014)''. \\
A02-13 & \sevM & 3.3 y Cuadro & El texto descarta métodos ``aplicados a \ldots\ ríos australianos, praderas canadienses, Corea del Sur'' sin citarlos, aunque están en \texttt{references.bib} (\texttt{andrew2023}, \texttt{barrow2019}, \texttt{ryu2024}). & Se agrega una fila al cuadro con esas tres citas. \\
A02-14 & \sevB & Bibliografía & La clave \texttt{andrew2023} usa el nombre de pila del primer autor (John, Andrew). & Clave \texttt{john2023} en la bibliografía revisada. \\
A02-15 & \sevB & Cuadro & Orden: HS2012 aparece antes que HS2009/2011, aunque estos originan la partición; Gopika figura como ``(EEP)'' cuando su dominio es 30°S--30°N. & Orden lógico; dominio corregido. \\
A02-16 & -- & 3.1 y Cuadro & Verificados sin cambios de fondo: Álvarez (título, CMIP5, TSM y nivel del mar); Keller (tendencias observadas 1974--2004 en Niño~3.4 en su Tabla~1; TOE de pCO$_2$ y pH $>30$ años frente a Perú y Chile; TOE de TSM de 45--90 años); Zhong; Lois; Poul (S/N con varianzas combinadas, umbral 2). & Se conservan, con el detalle verificado. \\
\end{hallazgos}

\section{Metadatos bibliográficos}

Los 22 DOI de \texttt{references.bib} resuelven en Crossref con autores, año, revista y volumen correctos. Cuatro rangos de páginas difieren del registro oficial y se corrigen en \archivo{revision_e1/informe/references_rev.bib}: zhong2023 (2535--2549), chadwick2019 (1635--1647), hawkins2009 (1095--1108) y szabo2016 (239--261). El detalle está en la auditoría A09.

\section{Extensión}

Original: Antecedentes 591 palabras y Literatura 381 (972 en total). Propuesto: 1055 
palabras. El aumento se debe a que la descripción de Hawkins--Sutton y Szabó ahora es correcta y completa, y a la fila nueva del cuadro.

\section{Texto propuesto}

\begin{propuesto}
% ============================================================
\section{Revisión científica}
\label{sec:marco}
% ============================================================

\subsection{Antecedentes}
\label{sec:antecedentes}

\textbf{Partición de incertidumbre (Hawkins y Sutton).} \citet{hawkins2009,hawkins2011precip} establecieron la partición de la incertidumbre de una proyección climática en tres fuentes: variabilidad interna ($V$, no reducible, propia de la dinámica caótica del sistema climático), incertidumbre de modelo ($M$, dispersión entre las respuestas de distintos modelos ante un mismo forzante) e incertidumbre de escenario ($S$, dispersión entre trayectorias de emisiones). Cada simulación se ajusta con un polinomio de cuarto orden, el residuo del ajuste representa la variabilidad interna y, al tratar las tres fuentes como independientes, la varianza total es su suma, $T = V + M + S$. \citet{hawkins2012} llevaron el marco al \emph{Time of Emergence} por modelo individual: la señal local se obtiene regresando la variable sobre la temperatura media global suavizada con el mismo polinomio, el ruido es la desviación estándar interanual del residuo y el TOE es el primer año en que la razón señal-ruido supera 1 o 2. Estimaron además que la incertidumbre del propio TOE es de al menos 30 años y llega a 60 en algunas regiones. La partición en tres fuentes es la base del análisis de incertidumbre previsto en el TdR para el Entregable~4.

\textbf{Aplicación de Szabó y Szépszó.} \citet{szabo2016} implementan el método de Hawkins y Sutton para temperatura y precipitación sobre Europa (15 modelos CMIP3, tres escenarios SRES, referencia 1971--2000) e introducen cuatro cambios: la variabilidad interna se recalcula cada año sobre los 30 años previos, en lugar de asumirse constante; se usan medias anuales y estacionales en vez de decadales; los modelos no se ponderan; y el TOE se define con la variabilidad interna como único ruido, mientras que la razón señal-ruido conserva la incertidumbre total. Esta formulación es la que el equipo técnico adopta como método M1 (Sección~\ref{sec:toe_metodos}).

\textbf{Trabajo de Bruno Ramírez.} \citet{bruno2023}, trabajo de suficiencia profesional de la Universidad Nacional Agraria La Molina, aplica la metodología de \citet{szabo2016} a temperatura y precipitación sobre Sudamérica con 53 modelos CMIP6, periodo de referencia 1981--2010 (1981--2014 en este proyecto) y los escenarios SSP2-4.5 y SSP5-8.5, calculando variabilidad interna, razón señal-ruido y TOE. Su validación de modelos combina un componente físico (posición e intensidad del anticiclón del Pacífico Sur, la Alta de Bolivia y la Baja Amazónica) y uno estadístico frente a ERA5 (correlación de Pearson, sesgo, RMSE y diagrama de Taylor). Es el antecedente metodológico más cercano: muestra que M1 ya se aplicó con CMIP6 en Sudamérica, aunque sobre variables atmosféricas continentales y no sobre las cajas oceánicas Niño~3.4 y Niño~1+2. Su lista de modelos se tomó como referencia para la selección de este entregable, que abarca todo el catálogo CMIP6 disponible (Sección~\ref{sec:desviacion}).

\textbf{Trabajo de Álvarez Sánchez.} \citet{alvarez2024}, también trabajo de suficiencia profesional de la misma universidad, proyectó la temperatura superficial del mar y el nivel del mar en las costas de Sudamérica y el Pacífico tropical con modelos CMIP5. Comparte con este proyecto el dominio oceánico y permite, en entregables posteriores, contrastar las generaciones CMIP5 y CMIP6.

\textbf{Niño costero.} \citet{takahashi2019} reconstruyen con observaciones in situ el Niño costero de 1925: calentamiento intenso en el Pacífico oriental lejano con condiciones frías en el Pacífico central, a diferencia de los eventos extraordinarios de 1982--1983 y 1997--1998. Las ondas Kelvin ecuatoriales tuvieron un papel menor en su inicio; el evento se asocia con vientos del norte a través del ecuador, con la intensificación de la zona de convergencia intertropical al sur del ecuador y con realimentaciones océano-atmósfera locales. Existen, por tanto, dos tipos de eventos de fuerte impacto en Niño~1+2, con dinámicas distintas, lo que justifica analizar esta caja por separado de Niño~3.4.

La revisión de metodologías de TOE que pide la actividad~(a) se extendió más allá de Hawkins-Sutton y Szabó; la Sección~\ref{sec:referencias} sintetiza la bibliografía revisada, de la que se deriva el método complementario ToD (Sección~\ref{sec:toe_metodos}).

\subsection{Literatura sobre \emph{Time of Emergence}}
\label{sec:referencias}

El Cuadro~\ref{tab:referencias} resume la bibliografía sobre TOE revisada para este proyecto: los antecedentes de la Sección~\ref{sec:antecedentes} y literatura complementaria que sustenta la discusión del Entregable~4.

\begin{longtable}{@{}L{3.6cm}L{3.2cm}L{\dimexpr\textwidth-3.6cm-3.2cm-4\tabcolsep\relax}@{}}
\caption{Bibliografía sobre \emph{Time of Emergence} revisada para este proyecto.}
\label{tab:referencias} \\
\toprule
\textbf{Referencia} & \textbf{Dominio} & \textbf{Aporte principal a este proyecto} \\
\midrule
\endfirsthead
\toprule
\textbf{Referencia} & \textbf{Dominio} & \textbf{Aporte principal a este proyecto} \\
\midrule
\endhead
\citet{hawkins2009,hawkins2011precip} & Global y regional (CMIP3) & Partición de la incertidumbre en tres fuentes ($T=V+M+S$) \\
\citet{hawkins2012} & Global, temperatura del aire (CMIP3) & TOE por modelo individual; umbrales S/N de 1 y 2 \\
\citet{szabo2016} & Europa (CMIP3) & Método M1: variabilidad interna móvil; TOE con ruido igual a $V$ \\
\citet{bruno2023} & Sudamérica, temperatura y precipitación (CMIP6) & Aplicación de M1 a CMIP6 en Sudamérica \\
\citet{alvarez2024} & Costas de Sudamérica y Pacífico tropical (CMIP5) & Antecedente institucional; comparación CMIP5 frente a CMIP6 \\
\citet{takahashi2019} & Pacífico oriental lejano & Dinámica del Niño costero \\
\citet{gopika2024} & Océanos tropicales, 30°S--30°N (CMIP5/6) & \emph{Time of Detection} (ToD), método complementario adoptado \\
\citet{zhong2023} & Pacífico tropical (CMIP6 y grandes ensambles) & Regímenes de incertidumbre distintos para el estado medio y para el ENOS \\
\citet{larson2025} & Océano global (CESM2, 1900--2014) & La variabilidad de la circulación forzada por viento retrasa el TOE de la TSM \\
\citet{poul2025} & Mar del Norte, Báltico y Atlántico Norte & Razón S/N con varianzas combinadas; la surgencia retrasa el TOE \\
\citet{keller2014} & Océano global, biogeoquímica (CMIP5) & TOE de TSM de 45--90 años; TOE largo frente a Perú y Chile \\
\citet{chadwick2019} & Chile semiárido y mediterráneo & TOE local por potencia estadística; relación entre CV y TOE \\
\citet{lois2026} & Océano Austral (CMIP6) & \emph{Storylines} y TOE; criterio alternativo de selección de modelos \\
\citet{giorgi2009} & Precipitación global (CMIP3) & TOE de \emph{hot-spots}; emergencia tardía en la cuenca amazónica \\
\citet{john2023,deng2022,barrow2019,ryu2024,ryu2025,im2021,nguyen2018} & Australia, Canadá, Corea del Sur, oeste de EE.\,UU., CORDEX-CORE, precipitación global & Otras aplicaciones de TOE revisadas y no adoptadas (hidrología, extremos de temperatura, estrés térmico, aridez, precipitación) \\
\bottomrule
\end{longtable}

El hallazgo transversal de esta bibliografía es que no existe un método estándar de TOE y que las decisiones metodológicas desplazan la fecha resultante en décadas. El umbral varía entre estudios (S/N de 1 en \citealp{szabo2016}, de 1 y 2 en \citealp{hawkins2012}, de 1{,}6 en \citealp{gopika2024} y de 2 en \citealp{poul2025}); \citet{hawkins2012} estiman una incertidumbre del TOE de 30 a 60 años, y \citet{chadwick2019} encuentran diferencias de más de 40 años según el percentil de tendencia de los modelos. El equipo técnico adopta el método de Szabó (M1) y añade el ToD de Gopika como segundo método independiente, para reportar un rango de fechas y no un único número.

\end{propuesto}

\bibliographystyle{plainnat}
\bibliography{../informe/references_rev}
\end{document}
